% !TEX TS-program = pdflatex
\documentclass[10pt,a4paper]{altacv}

\geometry{left=1.5cm,right=1.5cm,top=1.5cm,bottom=1.0cm,footskip=2\baselineskip}

\usepackage[T1]{fontenc}
\usepackage[utf8]{inputenc}
\usepackage[default]{lato}
\usepackage{microtype}
\DisableLigatures{encoding=*, family=*}

\definecolor{Navy}{HTML}{0A2540}
\definecolor{SlateGrey}{HTML}{2E2E2E}
\definecolor{LightGrey}{HTML}{666666}
\colorlet{heading}{Navy}
\colorlet{accent}{Navy}
\colorlet{emphasis}{SlateGrey}
\colorlet{body}{LightGrey}

\renewcommand{\itemmarker}{{\small\textbullet}}
\renewcommand{\ratingmarker}{\faCircle}

\renewcommand{\emailsymbol}{\faEnvelope}
\renewcommand{\phonesymbol}{\faPhone}
\renewcommand{\linkedinsymbol}{\faLinkedin}
\renewcommand{\locationsymbol}{\faMapMarker}
\makeatletter
\renewcommand{\printinfo}[2]{\mbox{\textcolor{accent}{\normalfont #1}\hspace{0.35em}#2}\hspace{1.1em}}
\newcommand{\objective}[1]{\def\@objective{#1}}
\newcommand{\profile}[1]{\def\@profile{#1}}
\makeatother
\graphicspath{{./}{../}{../../templates/}}
\usepackage[hidelinks]{hyperref}
\hypersetup{
  pdftitle={Curriculum Vitae -- Achraf Izelmadene},
  pdfauthor={Achraf Izelmadene},
  pdfsubject={Ingénieur Génie Électrique et Systèmes Intelligents}
}
\photo{2.8cm}{assets/imageee.png}

\renewcommand{\cvsection}[2][]{%
  \vspace{3pt}%
  {\color{heading}\LARGE\bfseries #2 \hspace{0.2em}\leaders\vrule height 4pt depth -2pt\hfill\mbox{}}\par\vspace{2pt}
}
\renewcommand{\divider}{\textcolor{body!30}{\hdashrule{\linewidth}{0.6pt}{0.5ex}}\vspace{1pt}}

\renewcommand{\cvevent}[4]{%
  {\large\bfseries\color{accent}#1\par}
  \vspace{1pt}\normalsize
  \ifstrequal{#2}{}{}{%
    \textbf{\color{emphasis}#2\ifstrequal{#4}{}{}{ - #4}}%
  }%
  \hfill
  \ifstrequal{#3}{}{}{{\small#3}}%
  \par\vspace{1pt}
}

\makeatletter
\renewcommand{\makecvheader}{%
  \begingroup
    \noindent
    \begin{minipage}[c]{\@photodiameter}
      \begin{tikzpicture}
        \useasboundingbox (-0.5\linewidth,-0.5\linewidth) rectangle (0.5\linewidth,0.5\linewidth);
        \clip (0,0) circle (0.5\linewidth);
        \node at (0,-0.1\linewidth) {\includegraphics[width=1.05\linewidth]{\@photo}};
      \end{tikzpicture}%
    \end{minipage}\hspace{1.2em}%
    \begin{minipage}[c]{\dimexpr\linewidth-\@photodiameter-1.2em}
      \raggedright\color{emphasis}%
      {\fontsize{22}{24}\selectfont\bfseries\MakeUppercase{\@name}\par}\vspace{2pt}
      {\large\bfseries\color{accent}\@tagline\par}\vspace{3pt}
      {\small\bfseries\color{emphasis}\@objective\par}\vspace{2pt}
      {\small\color{body}\@profile\par}\vspace{4pt}
      {\footnotesize\bfseries\color{emphasis}\@personalinfo\par}
    \end{minipage}\par
  \endgroup\smallskip
}
\makeatother

\begin{document}

\name{Izelmadene Achraf}
\tagline{Élève Ingénieur - Embedded Systems \& Edge AI}
\objective{À la recherche d'un stage d'ingénieur de 6 mois dès février 2027 : Concevoir et industrialiser des solutions d'Intelligence Artificielle sur systèmes embarqués contraints.}
\profile{Passionné par l'innovation, je m'investis dans l'intégration de technologies de pointe et la conception d'outils d'automatisation. Mes compétences en systèmes de contrôle et en analyse des performances répondent parfaitement aux enjeux de vos équipements industriels.}
\personalinfo{%
  \phone{+212699098736}%
  \email{\href{mailto:achraf.izelmadene@gmail.com}{achraf.izelmadene@gmail.com}}%
  \linkedin{\href{https://www.linkedin.com/in/achraf-izelmadene/}{achraf-izelmadene}}%
  \location{Marrakech, Maroc}%
}

\makecvheader

\cvsection{EXPÉRIENCE}
\cvevent{Stagiaire Bureau d'Études}{Core Engineering Solutions}{juil. - août 2026}{Berrchid, Casablanca}
\begin{itemize}
\item Dimensionnement d'un poste HTA/BT et intégration d'une installation solaire PV en autoconsommation.
\item Réalisation de schémas unifilaires, bilans de puissance et notes de calcul (AutoCAD, Caneco BT, PVsyst).
\end{itemize}
\cvevent{Stagiaire en Unité de Production}{Menara Prefa}{juil. 2024}{Marrakech}
\begin{itemize}
\item Migration d'un automate industriel vers Siemens S7 sous TIA Portal, suivi de maintenance préventive/corrective.
\item Diagnostic des pannes électriques et participation à l'amélioration continue des processus industriels.
\end{itemize}

\cvsection{PROJETS}
{\large\bfseries\color{accent} Vision AI pour Langue des Signes (Edge AI \& TensorFlow Lite)}\\[0.05cm]
\begin{itemize}
  \item Développement d'un pipeline d'inférence temps réel combinant Transformer, GRU et CNN (analyse spatio-temporelle).
  \item Optimisation des calculs et déploiement avec TensorFlow Lite pour exécution à très faible latence.
\end{itemize}

{\large\bfseries\color{accent} Détection d'anomalies Smartgrid (Séries temporelles)}\\[0.05cm]
\begin{itemize}
  \item Comparaison d'algorithmes (KMeans, DBSCAN) sur des flux de données IoT issus de réseaux électriques.
  \item Évaluation de métriques avancées et développement d'un dashboard de supervision et d'analyse des anomalies.
\end{itemize}

{\large\bfseries\color{accent} Precision Agriculture IoT Ecosystem (TinyML \& Microcontrôleurs)}\\[0.05cm]
\begin{itemize}
  \item Déploiement d'un réseau IoT radio (LoRa) reliant des nœuds capteurs (Arduino) à une unité décisionnelle (ESP32).
  \item Intégration de logique floue et Edge AI embarquée pour la gestion autonome de l'irrigation et détection.
\end{itemize}

\cvsection{FORMATION}
\cvevent{Étudiant en cycle d'ingénieur en Génie Électrique et Systèmes Intelligents}{École Marocaine des Sciences de l'Ingénieur}{2025 - Présent}{Marrakech}
\small Systèmes Embarqués, Machine Learning, Deep Learning, Traitement du signal, Capteurs Intelligents, Automatique,\\
\small Smart Grids, Réseaux Locaux Industriels, Robotique Industrielle, Supervision \& Digital Twin, Gestion de projet...

\cvevent{Licence SMP, option Automatique}{Faculté des Sciences Semlalia}{2020 - 2024}{Marrakech}
\small Électromagnétisme, Électronique, Traitement du signal, Automatique, Instrumentation, Informatique Industrielle...

\cvsection{COMPÉTENCES}
\begin{tabular}{@{}ll@{}}
\textbf{Programmation \& IA :} & Python, C/C++, TinyML, TensorFlow Lite, ONNX, Machine Learning, AutoML, MLOps \\[2pt]
\textbf{Systèmes Embarqués :} & ESP32, Arduino, Microcontrôleurs, Architectures IoT, MATLAB, Simulink \\[2pt]
\textbf{Compétences transversales :} & Analyse de données, Expérimentation, Résolution de problèmes, Autonomie \\
\end{tabular}

\cvsection{CERTIFICATIONS}
\begin{itemize}[leftmargin=*, label=\textbullet, itemsep=2pt]
  \item \textbf{Foundations of IoT Systems and Industrial Automation} \hfill {\color{emphasis}LearnQuest (Coursera)}
  \item \textbf{Sensors and Sensor Circuit Design} \hfill {\color{emphasis}University of Colorado Boulder (Coursera)}
  \item \textbf{Certified 3DEXPERIENCE Collaborative Industry Innovator} \hfill {\color{emphasis}Dassault Systèmes}
  \item \textbf{MATLAB / Simulink Onramp} \hfill {\color{emphasis}MathWorks}
\end{itemize}

\cvsection{ACTIVITÉS \& COMPÉTITIONS}
\begin{itemize}[leftmargin=*, label=\textbullet, itemsep=2pt]
  \item \textbf{Hack The Intelligence} (AISEC 2026) - \textbf{1\textsuperscript{er} Prix} : {\small Développement du projet Smart Agriculture IoT.}
  \item \textbf{Vision AI Competition} (EMSI 2026) - \textbf{2\textsuperscript{ème} Prix} : {\small Conception d'un système de traduction en langue des signes (ASL).}
  \item \textbf{Congrès COCIA'26} - \textbf{Présentation Scientifique} : {\small Conception du prototype d'éclairage piézoélectrique SmartStruct.}
  \item \textbf{Hackathon RallyIA} (2026) - \textbf{Participation} : {\small Conception de l'architecture du projet "FindSpot" (Parking Intelligent).}
\end{itemize}


\cvsection{LANGUES \& DIVERS}
\textbf{Arabe :} Langue maternelle \hfill \textbf{Français :} DELF B2 \hfill \textbf{Anglais :} EFSET C2\\[2pt]
\textbf{Permis de conduire :} B \hfill \textbf{Loisirs :} Lecture, Échecs, Électronique DIY


\end{document}
